\section{QLBS: Q-learning Black-Scholes Model}

\subsection{A Brief Review}

We quickly introduce the QLBS model proposed in \cite{halperin2020qlbs}.
Consider a sequence of asset prices $\left\{ S_{t}\right\} _{t=0,1,\dots}$,
adapted under the filtration $\left\{ \cF_{t}\right\} $, upon which
we wish to build an option with payoff function $h$ at the maturity
time $T$. The option is realized as a hedge portfolio which consists
of some shares $u_{t}$ of the underlying asset and the risk-free
deposit $B_{t}$. Let 
\[
\Pi_{t}:=u_{t}S_{t}+B_{t}
\]
denote the value of the portfolio at time $t$. To fulfill the option
contract, the holding position is cleared at the terminal time $T$
and is fully converted into cash position, i.e.
\[
\Pi_{T}=B_{T}=h\left(S_{T}\right).
\]
The deposit position at a particular time is solved via the self-financing
condition which requires that the instantaneous value of the portfolio
is kept same before and after the re-hedging operation:
\begin{equation}
u_{t}S_{t+1}+e^{r\Delta t}B_{t}=u_{t+1}S_{t+1}+B_{t+1}\label{eq:self-financing}
\end{equation}
where $r$ stands for the risk-free interest rate. Eqn. \ref{eq:self-financing}
can be used to solve how much money is needed to cover future trading
activities.

A state is defined as a pair of an integer and a real number $\left(t,X_{t}\right)$,
where 
\[
X_{t}=-\left(\mu-\frac{\sigma^{2}}{2}\right)t+\log S_{t}
\]
is the compensated logarithm price. The action space is the real number,
indicating how much shares are hedged at a particular time. A policy
is a mapping from the state space to the action space, i.e.
\[
\pi:\left(t,X_{t}\right)\mapsto a_{t}.
\]
Notice that we use $a_{t}$ for the log-processed input $X_{t}$ while
$u_{t}$ for the asset price $S_{t}$ in normal scale. The policy
may depend on other macro factors, e.g. interest rate $r$, volatility
$\sigma$, total maturity time $T$, the strike price $K$ if the
option is of call/put type.

The reward function is derived from the Bellman's optimality equation
where we define the value function in the first place. The idea is
to minimize the money needed to initiate the hedge portfolio as well
as to minimize the volatility throughout the trading periods. Given
a hedging strategy $\pi$, the value function is defined as
\begin{equation}
V_{t}^{\pi}\left(X_{t}\right)=\bE^{\pi}\left[-\Pi_{t}\left(X_{t}\right)-\lambda\sum_{\tau=t}^{T}e^{-r\left(\tau-t\right)}\text{Var}\left[\Pi_{\tau}\left(X_{\tau}\right)|\cF_{\tau}\right]|\cF_{t}\right]\label{eq:qlbs-value-function}
\end{equation}
where $\lambda$ is the risk aversion factor. The reward function
can be derived by matching the corresponding terms in the Bellman's
equation:
\[
R_{t}\left(X_{t},a_{t},X_{t+1}\right):=\gamma a_{t}\Delta S_{t}-\lambda\text{Var}\left[\Pi_{t}|\cF_{t}\right]
\]
where $\gamma:=e^{-r\Delta t}$ is the discounting factor. The connection
between the value function and option pricing is that the option price
is given by the minus optimal Q-function.

\subsection{Make QLBS Interactive}

Elegant as the vanilla QLBS approach, it is not a true RL problem
since the optimal policy $\pi^{*}$ is solved analytically without
any reinforcement learning techniques. In fact, the author derives
the Bellman's equation for the optimal Q-function from Eqn. \ref{eq:qlbs-value-function}
\begin{equation}
Q_{t}^{*}\left(X_{t},a_{t}\right)=\gamma\bE_{t}\left[Q_{t+1}\left(X_{t+1},a_{t+1}^{*}\right)+a_{t}\Delta S_{t}\right]-\lambda\gamma^{2}\bE_{t}\left[\widehat{\Pi}_{t+1}^{2}-2a_{t}\widehat{\Pi}_{t+1}\Delta\widehat{S}_{t}+a_{t}^{2}\left(\Delta\widehat{S}_{t}\right)^{2}\right]\label{eq:qlbs-optimal-Q-function}
\end{equation}
which admits the optimal policy in closed form since Eqn. \ref{eq:qlbs-optimal-Q-function}
is a quadratic function in $a_{t}$. Such direct approach is feasible
if provided with abundant data on the correlation structure of the
portofolio value $\widehat{\Pi}$ and the stock price change $\Delta\widehat{S}$,
but it fails to generalize beyond this simple setting. Besides, the
portofolio value process $\Pi_{t}$ is in general non-adapted due
to how the self-financing condition (Eqn. \ref{eq:self-financing})
works. 

To deal with the aforementioned issues, we propose a modified QLBS
model which is
\begin{enumerate}
\item fully adapted with respect to the given filtration $\left\{ \cF_{t}\right\} $,
\item compatible with the transaction cost proposed in Sec. \ref{subsec:Trading-Cost},
and
\item works well with value-based or policy-based learning algorithms.
\end{enumerate}
We start by modifying the value function. A first problem lies in
the fact that the reward are not homogeneous over time. Eqn. \ref{eq:qlbs-value-function}
assigns the (negative) cashflow with a risk part as the reward, but
the terminal step gets the option payoff which is in general much
larger than the previous steps. Technically speaking, an agent could
notice this heterogenity since the time $t$ aligns with the terminal
time $T$, but it is rather difficult in practice to figure out this
situation. Thus, we propose the modified value function 
\begin{equation}
V_{t}^{\pi}\left(X_{t}\right)=\bE_{t}^{\pi}\left[-\left(1-\frac{t}{T}\right)\Pi_{t}\left(X_{t}\right)-\lambda\sum_{\tau=t}^{T}\gamma^{\tau-t}\sqrt{\text{Var}\left[\Pi_{\tau}\left(X_{\tau}\right)\right]}\right].\label{eq:qlbs-modified-V}
\end{equation}
Our improvement is two folds:
\begin{enumerate}
\item We weight the portfolio term by a diminishing factor $\left(1-\frac{t}{T}\right)$.
This factor does not impact the starting estimate at $t=0$ and it
fully vanishes at $t=T$. We point out that introducing this factor
will break the temporal symmetry so that the intermediate estimate
does not correspond to option pricing from the intermediate time steps.
\item We take a square root of the variance terms so that they turn into
standard deviation. This helps to keep the value estimate dimensionless
and robust.
\end{enumerate}
The reward function changes accordingly to
\begin{equation}
R_{t+1}\left(X_{t},a_{t}\right)=\bE_{t}^{\pi}\left[-\left(1-\frac{t}{T}\right)\Pi_{t}\left(X_{t}\right)+\left(1-\frac{t+1}{T}\right)\Pi_{t}\left(X_{t}\right)-\lambda\sqrt{\text{Var}\left[\Pi_{t}\left(X_{t}\right)\right]}\right]\label{eq:qlbs-modified-reward}
\end{equation}
with action at time $t$ to be $a_{t}$ and remaining actions following
the current policy $\pi$. Considering the effect of transaction costs,
the portfolio value process $\Pi_{t}$ is calculated backwards from
the modified self-financing condition
\begin{equation}
e^{r\Delta t}\left(\Pi_{t}-u_{t}S_{t}\right)+u_{t}S_{t+1}=\Pi_{t+1}+\text{TC}\left(u_{t+1}-u_{t},S_{t+1}\right)\label{eq:qlbs-modified-self-financing}
\end{equation}
with terminal condition $\Pi_{T}=h\left(S_{T}\right)$ unchanged.
We point out that we wrap the cashflow part with an conditional expectation
at time $t$ so that we don't run into adaptedness issues. We briefly
illustrate the modified QLBS model in Fig. \ref{fig:qlbs-illustration}.

\begin{figure}[]
\begin{centering}
\includegraphics[width=0.6\linewidth]{backup/qlbs-illustration}
\par\end{centering}
\caption{A brief illustration on the modified QLBS environment.}

\label{fig:qlbs-illustration}
\end{figure}


\subsection{Specification}

In this section, we describe how to set-up the QLBS environment and
the necessary numerical procedures.

\subsubsection{Environment \label{subsec:Environment}}

The environment is responsible for keeping track of the asset price
and portfolio value based on the actions provided by the agent. With
a given set of parameters $r,\mu,\sigma,T$, the asset prices are
a set of geometrical brownian motion paths $\left\{ S_{t}\right\} $
which solves
\[
\d{S_{t}}=\mu S_{t}\d t+\sigma S_{t}\d{W_{t}}
\]
where $W_{t}$ refers to the standard brownian motion according to
the filtration $\left\{ \cF_{t}\right\} $. At each time step $t$,
the normalized price $X_{t}=-\left(\mu-\frac{\sigma^{2}}{2}\right)t+\log S_{t}$
and $t$ are provided to the agent, waiting for the response of the
hedge position $a_{t}$. Then, the reward $R_{t+1}$, as a conditional
expectation specified in Eqn. \ref{eq:qlbs-modified-reward}, is computed
empirically by averaging samples from a fixed number of additional
trajectories under the current policy. At the beginning of each episode,
the parameters are adjusted in a random fashion to help the agent
explore different settings and help avoid overfitting; the adjustment
obeys a Poisson process with intensity $\varUpsilon$.

\subsubsection{Agent Parametrization \label{subsec:Agent-Parametrization}}

The agent is fully responsible for determining the hedge position
$a_{t}$ under a given normalized price at each time step. The policy
$\pi$, whether stochastic or deterministic, depends on these input
variables as well as the environment parameters
\[
a_{t}\sim\pi\left(X_{t},t;r,\mu,\sigma,T,K,\lambda\right).
\]
In pratice, we prefer a stochastic policy since it encourages exploration
which is helpful to excape local minima. To further simplify the sampling
procedure, we restrict our policy spaces to Gaussian distribution
where the agent determines the mean and standard deviation, i.e.
\[
\pi=\cN\left(\mu_{\pi},\sigma_{\pi}\right)
\]
(where the subscripts $\pi$ are used to distinguish these parameters
from the environmental ones). The statistics $\mu_{\pi}$ and $\sigma_{\pi}$
are parametrized by two separate neural networks with the Resnet \cite{he2016deep}
skip-connection structure. The Resnet structure is composed of three
parts:
\begin{enumerate}
\item Pre-processing $T_{\text{lift}}\left(x\right):=\Xi\left(w_{\text{lift}}^{T}x+b_{\text{lift}}\right)$,
that lifts the (8-dimensional) input to the latent dimension by an
affine transform and an activation funciton $\Xi$;
\item Chain of transforms $\left\{ T^{\left(k\right)}\right\} $ in a fixed-point
iteration style, with each transform combining the identity functions
and a series of alternating affine transforms $\left\{ Z_{l}^{\left(k\right)}\right\} $
and activation $\Xi$, i.e.
\[
T^{\left(k\right)}:=\Xi\circ\left[\text{id}+Z_{n_{k}}^{\left(k\right)}\circ\Xi\circ Z_{n_{k}-1}^{\left(k\right)}\circ\cdots\circ\Xi\circ Z_{1}^{\left(k\right)}\right]
\]
\item Post-processing $T_{\text{project}}\left(x\right):=w_{\text{project}}^{T}x+b_{\text{project}}$,
that projects the latent representations onto the target space.
\end{enumerate}
Thus, the Resnet realization can be formally written as $T_{\text{project}}\circ T^{\left(k\right)}\circ T^{\left(k-1\right)}\circ\cdots\circ T^{\left(1\right)}\circ T_{\text{lift}}$.
We include an illustration in Fig. \ref{fig:qlbs-resnet}.

\begin{figure}[]
\begin{centering}
\includegraphics[width=0.8\linewidth]{backup/resnet-illustration}
\par\end{centering}
\caption{An illustration on a Resnet-motivated network structure.}

\label{fig:qlbs-resnet}
\end{figure}


\subsubsection{Learning Algorithm}

Since the action space $a_{t}\in\bR$ is a continuous space, it is
nature to adopt a policy-based method. Here we opt-in the classical
REINFORCE algorithm with baseline where the policy components relies
on the value estimator to learn quickly and reliably while the value
estimator learns from empirical averages of Monte Carlo samples. The
value estimator, parameterized by a neural network, uses the same
Resnet structure as mentioned in the previous section. We refer to
\cite{sutton2018reinforcement} for more details.

In our implementation, the policy network and the value network have
the same latent dimension 10 and they are composed of two Resnet blocks
with two hidden affine transforms. We use the Adam optimizers to update
the networks, one for each, and the learning rate is set to $10^{-4}$
by default.

\subsection{Experiments and Results}

We include a variety of experiments, starting from a demonstration
that shows the agent learns over time, moving to a comparison with
the Black-Scholes baseline model, and finally exploring other directions
that have a stronger connection to finance.

\subsubsection{Demonstration on learning\label{subsec:Demonstration-on-learning}}

We start by showing that the policy gradient algorithm has been correctly
implemented and the agent learns well over time. The environment parameter
is set to $r=0.01,\sigma=0.1,\mu=0,T=5,K=1,S_{0}=1,\Delta t=1,\lambda\in\left\{ 0,1,2,3\right\} ,\epsilon=0$
in Sec. \ref{subsec:Demonstration-on-learning}, \ref{subsec:Influence-of-risk},
and \ref{subsec:Effect-of-transaction} unless otherwise specified.
As shown in Fig. \ref{fig:qlbs-exp1-learning}, the episodic return
flatterns after training for around 3000 steps. It is worth noticing
that the cashflow return part 
\[
\sum_{t=0}^{T-1}\bE^{\pi}\left[-\left(1-\frac{t}{T}\right)\Pi_{t}\left(X_{t}\right)+\left(1-\frac{t+1}{T}\right)\Pi_{t}\left(X_{t}\right)\right]=\bE^{\pi}\left[\Pi_{0}\left(X_{0}\right)\right]
\]
is slightly smaller for a large risk parameter $\lambda$, implying
that the agent learns to make a trade-off between the cash-flow and
the risk component. 

To prepare for learning in a much broader setting, we allow the initial
price to be adjusted at a given intensity $\varUpsilon$. We mark
out the episodes that the adjustments take place (Fig. \ref{fig:qlbs-exp1-mutate})
and it seems that the agent can swiftly adapt to the change.


\begin{figure}[]
\caption{No adjustment on environment parameters. The shade indicates 95\%
confidence interval. Learning rate set to $10^{-4}$.}
\includegraphics[width=0.45\linewidth]{qlbs/experiment1/learning-curve.png}

\label{fig:qlbs-exp1-learning}

\label{fig:qlbs-exp1}
\end{figure}

\begin{figure}[]
\caption{Adjustment parameter (on the initial price) $\varUpsilon=0.005$,
indicated by purple lines. The shade indicates 95\% confidence interval.
Learning rate set to $10^{-4}$. }
\includegraphics[width=0.45\linewidth]{qlbs/experiment1/mutate.png}

\label{fig:qlbs-exp1-mutate}
\end{figure}


\subsubsection{Influence of risk parameter $\lambda$\label{subsec:Influence-of-risk}}

We wish to compare the optimal price under different choices of the
risk aversion parameter $\lambda$. A priori estimate is that a larger
$\lambda$ leads to a higher price learnt, since it penalizes improper
hedges harder. To complete the comparision, we also introduce the
vanilla BS hedging strategy $\pi_{\text{BS}}$ and learns the episodic
return using the same neural network parametrization. The results
are shown in Fig. \ref{fig:qlbs-exp2}. We remind the readers that
the option prices are defined as the negative value function, so a
lower option price implies a larger return learnt. As we expect, a
large $\lambda$ does lead to a higher price, both under BS policy
and under the learnt policy. In general, the policy learnt by the
neural network parametrization is a bit sub-optimal compared to the
BS policy, but it performs better when $\lambda=0.$

\begin{figure}[]
\begin{centering}
\includegraphics[width=0.75\linewidth]{qlbs/experiment2/option-price.png}
\par\end{centering}
\caption{Negative episodic returns as option prices.}

\label{fig:qlbs-exp2}
\end{figure}


\subsubsection{Effect of transaction cost\label{subsec:Effect-of-transaction}}

Next, we'd like to examine the effect of transaction costs. As indicated
by the self-financing condition (Eqn. \ref{eq:qlbs-modified-self-financing}),
a higher trading cost can significantly bias the portfolio value process
and consequently increase the optimal price learnt. As shown in Fig.
\ref{fig:qlbs-exp3}, we compare the price and hedge position learnt
under different friction parameters $\epsilon$. In general, a larger
$\epsilon$ does lead to a higher price learned (shown in the left
half) and the hedge position is usually always higher (in the right
half). 

\begin{figure}[]
\begin{centering}
\includegraphics[width=1\linewidth]{qlbs/experiment3/price-hedge.png}
\par\end{centering}
\caption{Option priced and hedged under different transaction cost parameters
$\epsilon$. $\lambda$ is set to 0.5 to introduce some risk concerns.}

\label{fig:qlbs-exp3}
\end{figure}


\subsubsection{Generalization power}

As a concluding setting, we wish to parametrize the policy and baseline
networks on not only the states but also the environment parameters
and examine the generalization power. We pick two set of parameters:
$r=0.01,\mu=0,\sigma=0.1,\lambda=0.5,\epsilon=0$ under the first
condition while $r=0.02,\mu=0.1,\sigma=0.2,\lambda=1.5,\epsilon=0.1$
under the second condition. We train a policy and baseline under these
two conditions mixed, i.e. they have the same probability of being
presented, with the switching process being a Poisson process. After
training for a while, we refine the agent to learn a third condition
where the parameters are set to the arithmic average of the given
two conditions. As shown in Fig. \ref{fig:qlbs-exp4}, we examine
the hedging strategy of the agent before and after fine tuning to
the third average condition. The result, that the fine tuning does
a great improvement, is not surprising since we don't expect a lot
of generalization power. However, one could expect such generalization
given enough amount of training time.

\begin{figure}[]
\begin{centering}
\includegraphics[width=1\linewidth]{qlbs/experiment4/price-hedge.png}
\par\end{centering}
\caption{Option priced and hedged under different conditions.}

\label{fig:qlbs-exp4}
\end{figure}